\documentclass[11pt,a4paper]{ltjsarticle}
\usepackage{amsmath,amssymb,amsthm,mathtools,bm,mathrsfs}
\usepackage{luatexja-fontspec}
\usepackage{enumitem,booktabs,longtable,array,needspace}
\usepackage{xcolor}
\usepackage{etoolbox,xr-hyper}
\usepackage[colorlinks=true,linkcolor=blue!45!black,citecolor=blue!45!black,urlcolor=blue!45!black]{hyperref}
\DeclareMathSizes{9}{9}{6}{5}
\DeclareMathSizes{9.85501}{9}{6}{5}
\IfFileExists{/usr/share/fonts/opentype/noto/NotoSerifCJK-Regular.ttc}{%
  \setmainjfont[Path=/usr/share/fonts/opentype/noto/]{NotoSerifCJK-Regular.ttc}%
}{\setmainjfont{HaranoAjiMincho-Regular}}
\setlength{\textwidth}{15.6cm}
\setlength{\oddsidemargin}{0.25cm}
\setlength{\evensidemargin}{0.25cm}
\setlength{\textheight}{22.6cm}
\setlength{\topmargin}{-0.2cm}
\setlength{\parindent}{1em}
\setlength{\parskip}{1pt}
\numberwithin{equation}{section}
\newtheorem{theorem}{定理}[section]
\newtheorem{maintheorem}{主定理}
\renewcommand{\themaintheorem}{\Alph{maintheorem}}
\newtheorem{lemma}[theorem]{補題}
\newtheorem{proposition}[theorem]{命題}
\newtheorem{corollary}[theorem]{系}
\newtheorem{remark}[theorem]{注意}
\newtheorem{definition}[theorem]{定義}
\newtheorem{example}[theorem]{例}
\newtheorem{conjecture}[theorem]{予想}
\renewcommand{\proofname}{証明}
% 三稿で共通に用いる記号。本文と証明はそれぞれ一つの稿に置く。
\providecommand{\ZZ}{\mathbb Z}
\providecommand{\cK}{\mathcal K}
\providecommand{\cP}{\mathcal P}
\providecommand{\cD}{\mathcal D}
\providecommand{\cC}{\mathcal C}
\providecommand{\cE}{\mathcal E}
\providecommand{\cF}{\mathcal F}
\providecommand{\cM}{\mathcal M}
\providecommand{\fm}{\mathfrak m}
\providecommand{\Perf}{\operatorname{Perf}}
\providecommand{\Ku}{\operatorname{Ku}}
\providecommand{\MF}{\operatorname{MF}}
\providecommand{\HH}{\operatorname{HH}}
\providecommand{\Dsg}{D_{\mathrm{sg}}}
\providecommand{\RHom}{\operatorname{RHom}}
\providecommand{\REnd}{\operatorname{REnd}}
\providecommand{\Ind}{\operatorname{Ind}}
\providecommand{\Def}{\operatorname{Def}}
\providecommand{\thick}{\operatorname{thick}}
\providecommand{\Spf}{\operatorname{Spf}}
\providecommand{\Id}{\operatorname{Id}}
\providecommand{\sym}{\operatorname{sym}}
\providecommand{\Cone}{\operatorname{Cone}}
\providecommand{\Pic}{\operatorname{Pic}}
\providecommand{\otL}{\mathbin{\otimes^{\mathbf L}}}
\newtheorem{problem}[theorem]{問題}

\newcommand{\C}{\mathbb C}
\newcommand{\CC}{\mathbb C}
\newcommand{\PP}{\mathbb P}
\newcommand{\OO}{\mathcal O}
\newcommand{\Sym}{\operatorname{Sym}}
\newcommand{\Spec}{\operatorname{Spec}}
\newcommand{\Proj}{\operatorname{Proj}}
\newcommand{\Hom}{\operatorname{Hom}}
\newcommand{\End}{\operatorname{End}}
\newcommand{\Ext}{\operatorname{Ext}}
\newcommand{\Aut}{\operatorname{Aut}}
\newcommand{\GL}{\operatorname{GL}}
\newcommand{\PGL}{\operatorname{PGL}}
\newcommand{\Gr}{\operatorname{Gr}}
\newcommand{\rank}{\operatorname{rank}}
\newcommand{\qgr}{\operatorname{qgr}}
\newcommand{\QGr}{\operatorname{QGr}}

% 他稿の番号には R (可換復元), M (モジュライ・降下), S (混合族) を添える。
% clean build では scripts/build_papers.py が三稿の aux を順に生成・収束させる。
\newcommand{\SetPaper}[1]{%
  \def\CurrentPaper{#1}%
  \ifstrequal{#1}{main}{}{\externaldocument[main-][nocite]{build/nc-cubic}[nc-cubic.pdf]}%
  \ifstrequal{#1}{moduli}{}{\externaldocument[moduli-][nocite]{build/cubic-moduli-descent}[cubic-moduli-descent.pdf]}%
  \ifstrequal{#1}{mixed}{}{\externaldocument[mixed-][nocite]{build/mixed-sklyanin-cubics}[mixed-sklyanin-cubics.pdf]}%
}
\makeatletter
\def\PaperTag@main{R}
\def\PaperTag@moduli{M}
\def\PaperTag@mixed{S}
\newcommand{\PaperRefText}[2]{\textnormal{[\csname PaperTag@#1\endcsname,~\ref{#1-#2}]}}
\newcommand{\PaperEqRefText}[2]{\textnormal{[\csname PaperTag@#1\endcsname,~\eqref{#1-#2}]}}
\DeclareRobustCommand{\paperref}[2]{\ifmmode\text{\PaperRefText{#1}{#2}}\else\PaperRefText{#1}{#2}\fi}
\DeclareRobustCommand{\papereqref}[2]{\ifmmode\text{\PaperEqRefText{#1}{#2}}\else\PaperEqRefText{#1}{#2}\fi}
\DeclareRobustCommand{\paperpageref}[2]{\textnormal{[\csname PaperTag@#1\endcsname,~p.~\pageref{#1-#2}]}}
\newcommand{\PaperReferenceGuide}{%
  本稿と関連する二稿の参照では、R を可換復元稿\cite{NCCubicRealization}、
  M をモジュライ・降下稿\cite{NCCubicModuliDescent}、
  S を混合族稿\cite{NCCubicMixedSklyanin}の略号とする。
  他稿の結果は、例えば定理 [R, 2.1] のように稿の略号とその稿内の番号を併記する。
  略号のない番号は本稿内の参照である。
}
\makeatother
